\documentclass[12pt]{article}
\usepackage{amsfonts,amsthm,amsmath,amssymb,mathtools}
\usepackage{mathrsfs}
\usepackage{enumitem}
\usepackage{tikz-cd}
\usepackage{geometry}
\usepackage{mlmodern}
\usepackage{graphicx}

\geometry{letterpaper, portrait, margin=1in}

\setcounter{section}{0}

\renewcommand{\thesection}{\S\arabic{section}}
\renewcommand{\thesubsection}{\arabic{section}}
\DeclareMathOperator{\im}{im}
\DeclareMathOperator{\id}{id}

\newtheorem{thm}{Theorem}
\newenvironment{theorem}[1]{
	\renewcommand{\thethm}{#1}
	\begin{thm}
		}{\end{thm}}

\newtheorem{lma}{Lemma}
\newenvironment{lemma}[1]{
	\renewcommand{\thelma}{#1}
	\begin{lma}
		}{\end{lma}}

\theoremstyle{definition}
\newtheorem{ex}{Assignment}
\newenvironment{exercise}[1]{
	\renewcommand{\theex}{#1}
	\begin{ex}
		}{\end{ex}}

\title{MTH 732 Topology - Homework 9}
\author{Connor Mooneyhan}
\date{April 2, 2026}

\begin{document}

\maketitle

\begin{ex}
	For a CW complex $X$, show that $(X^n, X^{n-1})$ is a good pair.
\end{ex}
\begin{proof}
	First, we show that $X^{n-1}$ is closed.
	If $X^n = X^{n-1}$, then $X^{n-1}$ is trivially closed in $X^n$.
	Otherwise, let $x \in X^n \setminus X^{n-1}$.
	Then $x$ must be in some $n$-cell $e^m$, which itself is a neighborhood of $x$ that does not intersect $X^{n-1}$.
	Thus $X^{n-1}$ is closed since its complement is open.

	Now we find a neighborhood which deformation retracts to $X^{n-1}$.
	If $X^n = X^{n-1}$, then $X^n$ is open and trivially deformation retracts to $X^{n-1}$, so we are done.
	Otherwise, consider $S = X^n \cap X^{n-1}$.
	Each point $s \in S$ is in the image of a boundary map of some $D^n_s$ (if there are multiple possibilities for $D^n_s$, pick one) with attaching map $g_s$.
	We can also define $U_s$, a neighborhood of $g_s^{-1}(s)$ in $D^n_s$ by taking the subset of $D^n_s$ where the radius is greater than 0.9.
	Let $p : X^{n-1} \bigsqcup D^n_\alpha \to X^{n}$ be the standard projection onto the quotient.
	Then we can define \begin{equation*}
		U = X_{n-1} \bigcup_{s \in S} p(U_s).
	\end{equation*}
	$U$ is open since $U \cap X^{n-1} = X^{n-1}$ which is open in $X^{n-1}$, and $U \cap e^n_\alpha$ is open for each $\alpha$ in the indexing set for $n$-cells by construction.
	Also by construction, $U$ deformation retracts onto $X^{n-1}$, so we conclude that $(X^n, X^{n-1})$ is a good pair.
\end{proof}

\begin{ex}
	Describe a CW complex structure on the torus $T^2$ having exactly one 0-cell, two 1-cells, and one 2-cell.
	Describe all attaching maps as well.
\end{ex}
\begin{proof}
	The fundamental polygon (a square) works great for this.
	We don't even need to subdivide it this time.
	We just have a 0-cell corresponding to the ``four corners,'' two 1-cells corresponding to the two pairs of opposite-identified sides, and the 2-cell filling it all in in the middle.
\end{proof}

\begin{ex}
	Let $v, e, f$ be positive integers.
	If $v - e + f = 2$, show that it is possible to construct a CW complex $X$ having exactly $v, e, f$ many cells in dimensions 0, 1, and 2, respectively, such that $X$ is homeomorphic to the 2-sphere. \textit{Hint: it might be easier to split this problem into three cases, according to whether $e \geq v - 1$, $e = v$, or $e < v$.}
\end{ex}
\begin{proof}
	Note that many of our constructions may be constructed as being ``on the sphere'', but this is just to give a sense of structure and should not affect the argument if constructed properly.

	If $e = v$, then $f = 2$.
	Since $v$ cannot be 0, consider first when $e = v = 1$.
	Then we have one vertex with a loop on it, and we can attach two copies of $D^2$ to that loop, thinking of them as being on ``both sides of it''.
	This is clearly a 2-sphere.

	Now if $v \geq 2$ (but still $e = v$), we first put vertices $N$ and $S$ at antipodal points, which will act as our north and south poles.
	We can then attach two edges, each going from $N$ to $S$, and then choose one of the edges to subdivide by adding $v-2$ vertices to it.
	Each time we do this, $e$ goes up by one as well, so the equality is maintained.
	Finally, we can attach our two copies of $D^2$ to the ``two'' original edges even though one has now been subdivided, since that will not affect the attaching maps.

	Now in the case where $e > v$ we follow the same construction as the previous paragraph until we have exhaused all the vertices we have available.
	We even attach two faces in the same way.
  Observe that we have $e - v$ more 1-cells to allocate and $f - 2$ more 2-cells (noting that since $v - e < 0$, the equation forces $f \geq 3$)\dots
  I'm not sure how to finish this case.

  Finally, consider where $e < v$.
  Then $v - e > 0$, so $f$ must be 1 since it is positive, which makes $v - e = 1$, and thus $v = e + 1$.
  In this case, we put two antipodal vertices ``on the sphere,'' and one edge connecting them.
  We subdivide this edge with new vertices (and thus new edges) as before, $v - 2$ times.
  We have now exhausted all the vertices and edges, so we attach our single face to the ``big edge'' from $N$ to $S$, wrapping it around so that, for instance, $[0, \pi]$ gets mapped from $N$ to $S$, and then $[\pi, 2\pi]$ gets mapped from $S$ to $N$ along the same ``big edge.''
  This is a 2-sphere, so we are done.
\end{proof}

\end{document}
